\chapter{System Overview of BED}

The Behavioral Engine for Demand (BED) is a closed-loop incentive system that
integrates behavioral modeling, constrained optimization, simulation, and
real-time experimentation. This chapter provides a high-level overview of the
BED architecture, including its conceptual components, data flow, and
operational logic. The goal is to establish a unified framework that connects
behavioral theory with practical incentive deployment.

\section{Conceptual Architecture}

BED is designed as a modular system composed of four core components:

\begin{enumerate}
    \item \textbf{Behavioral State Engine}: Tracks latent behavioral variables
    such as recency, habit, and responsiveness.
    \item \textbf{Hazard Model}: Predicts the probability of an event (e.g.,
    purchase, engagement) given the current behavioral state.
    \item \textbf{CMDP Policy Engine}: Allocates incentives optimally under
    budget and fairness constraints.
    \item \textbf{Experimentation and Logging Layer}: Captures real-time events
    and updates behavioral states based on observed outcomes.
\end{enumerate}

These components form a closed-loop system in which incentives influence
behavior, behavior influences state transitions, and state transitions inform
future incentives.

\section{Closed-Loop Data Flow}

The BED system operates through a continuous feedback loop:

\begin{enumerate}
    \item \textbf{Observation}: The system observes the user's current
    behavioral state.
    \item \textbf{Prediction}: The hazard model estimates the probability of an
    event occurring.
    \item \textbf{Decision}: The CMDP policy engine selects an incentive action
    subject to budget constraints.
    \item \textbf{Deployment}: The incentive is delivered to the user.
    \item \textbf{Outcome}: The user either engages (event) or does not engage
    (no event).
    \item \textbf{Update}: The behavioral state engine updates recency, habit,
    and responsiveness based on the outcome.
\end{enumerate}

This loop repeats for every user interaction, enabling adaptive and personalized
incentive strategies.

\section{Behavioral State Representation}

BED models behavior using three latent variables:

\begin{itemize}
    \item \textbf{Recency ($r$)}: Measures how recently the user engaged.
    \item \textbf{Habit ($h$)}: Captures the strength of repeated behavior.
    \item \textbf{Responsiveness ($\rho$)}: Represents sensitivity to incentives.
\end{itemize}

These variables form a Markovian state that evolves according to interpretable
update rules. Incentives influence state transitions by increasing the
likelihood of engagement and accelerating habit formation.

\section{Hazard-Based Event Prediction}

The hazard model estimates the probability of an event given the behavioral
state. BED uses a logistic hazard function:



\[
P(\text{event} \mid r, h, \rho) = \sigma(\beta_0 + \beta_r r + \beta_h h + \beta_\rho \rho)
\]



where $\sigma$ is the logistic function. This formulation allows BED to capture
nonlinear behavioral effects and heterogeneity in responsiveness.

\section{CMDP Policy Engine}

The incentive allocation problem is formulated as a Constrained Markov Decision
Process (CMDP). The CMDP ensures that:

\begin{itemize}
    \item incentives are allocated efficiently,
    \item long-term behavioral impact is maximized,
    \item budget and fairness constraints are respected.
\end{itemize}

The CMDP policy engine selects actions that balance immediate engagement with
long-term habit formation.

\section{Experimentation and Logging Layer}

To support real-world deployment, BED includes a lightweight experimentation
layer that logs events, treatments, and outcomes. This layer enables:

\begin{itemize}
    \item real-time state updates,
    \item A/B and multi-arm experiments,
    \item evaluation of policy performance,
    \item safe and privacy-aware data collection.
\end{itemize}

In BED v3, this layer is implemented using a serverless QR-to-cloud prototype
that logs synthetic events to a Firestore database.

\section{Summary}

The BED system integrates behavioral modeling, hazard-based prediction,
constrained optimization, and real-time experimentation into a unified
closed-loop architecture. This chapter establishes the conceptual foundation for
the behavioral model (Chapter 4), CMDP formulation (Chapter 5), and simulation
environment (Chapter 6) that follow.
